\begin{longtable}{c c p{9cm}}
\toprule
\textbf{Objetivo} & \textbf{Req. ID} & \textbf{Descripción del Requisito} \\
\midrule
\endfirsthead

% \caption[]{Requisitos funcionales detallados de la librería QTensor (continuación)} \\
\toprule
\textbf{Objetivo} & \textbf{Req. ID} & \textbf{Descripción del Requisito} \\
\midrule
\endhead
\bottomrule
\caption{Requisitos funcionales detallados de la librería QTensor.}
\label{tab:rf_detailed}
\endfoot

% --- Objetivo 1: Diseño e Implementación de la Librería ---
\multirow{5}{*}{\parbox{1.5cm}{\centering 1. Diseño y API}}
    & RF1.1 & El sistema permitirá definir una estrategia de cuantización global o por defecto para un modelo. \\
    \cmidrule{2-3}
    & RF1.2 & El sistema permitirá anular la estrategia global y asignar esquemas de cuantización específicos a capas individuales del modelo. \\
    \cmidrule{2-3}
    & RF1.3 & El sistema permitirá una granularidad a nivel de atributo, posibilitando la asignación de cuantizadores distintos para cada parámetro del modelo. En particular, para cualquier peso (e.g. kernel, bias de capas densas) y cualquier activación de cada capa, incluyendo capas con parámetros anidados (ej. RNN). \\
    \cmidrule{2-3}
    & RF1.4 & El sistema será compatible con la API de modelos \texttt{Sequential} y \texttt{Functional} de Keras. \\
    \cmidrule{2-3}
    & RF1.5 & El sistema soportará la serialización y deserialización de modelos cuantizados (formato .keras o SavedModel), preservando el estado completo de los cuantizadores, incluyendo sus parámetros entrenables. \\
\midrule

% --- Objetivo 2: Implementación de Cuantización Uniforme ---
\multirow{3}{*}{\parbox{1.5cm}{\centering 2. Cuant. Uniforme}}
    & RF2.1 & El sistema incluirá una clase \texttt{UniformQuantizer} que implementará la cuantización uniforme. Será configurable en número de bits y en el modo de operación (con signo o sin signo). \\
    \cmidrule{2-3}
    & RF2.2 & La clase \texttt{UniformQuantizer} expondrá su rango de cuantización (parámetro $\alpha$) como una variable entrenable, permitiendo al modelo optimizar dicho rango durante el entrenamiento. \\
    \cmidrule{2-3}
    & RF2.3 & La clase \texttt{UniformQuantizer} implementará un gradiente personalizado (STE) para permitir la propagación de gradientes, tanto para las entradas como para el parámetro entrenable $\alpha$. \\
    \cmidrule{2-3}
    & RF2.4 & La clase \texttt{UniformQuantizer} permitirá la inicialización del parámetro $\alpha$ mediante un valor fijo. \\
\midrule

% --- Objetivo 3: Implementación de Cuantización Flexible ---
\multirow{3}{*}{\parbox{1.5cm}{\centering 3. Cuant. Flexible}}
    & RF3.1 & El sistema incluirá una clase \texttt{FlexQuantizer} para la cuantización flexible. Será configurable en número de bits, número de niveles y modo de operación (con signo o sin signo). \\
    \cmidrule{2-3}
    & RF3.2 & La clase \texttt{FlexQuantizer} expondrá su rango ($\alpha$), los valores de los niveles de cuantización ($l$) y los umbrales de decisión ($t$) como variables entrenables. \\
    \cmidrule{2-3}
    & RF3.3 & La clase \texttt{FlexQuantizer} implementará un gradiente personalizado que permita el entrenamiento de los parámetros del modelo y del cuantizador (rango, niveles y umbrales). \\
\midrule

% --- Objetivo 4: Herramientas de Evaluación y Experimentación ---
\multirow{4}{*}{\parbox{1.5cm}{\centering 4. Evaluación}}
    & RF4.1 & La librería proveerá un `Callback` de Keras diseñado para recolectar y registrar la evolución de los parámetros entrenables de todos los cuantizadores del modelo en cada época de entrenamiento. \\
    \cmidrule{2-3}
    & RF4.2 & Se proporcionarán utilidades para visualizar los datos recolectados, permitiendo generar gráficos del estado y la evolución de los parámetros de los cuantizadores. \\
    \cmidrule{2-3}
    & RF4.3 & La librería incluirá una función para calcular la complejidad espacial teórica. Para cuantizadores uniformes, se basará en el número de bits fijos. Para cuantizadores flexibles, se basará en la entropía de los pesos cuantizados para simular una codificación de Huffman. \\
    \cmidrule{2-3}
    & RF4.4 & Se desarrollará al menos un script de ejemplo que demuestre el uso de la librería, permitiendo configurar el modelo, el dataset y la estrategia de cuantización. \\
\midrule

% --- Objetivo 5: Orquestación Modular de Experimentos ---
\multirow{3}{*}{\parbox{1.5cm}{\centering 5. Orquestación de experimentos}}
    & RF5.1 & El sistema definirá un sistema para desarrollar experimentos de comparación entre distintas estrategias de cuantización, permitiendo ejecutar flujos de trabajo completos (pre-entrenamiento, evaluación, cuantización, QAT, evaluación final) de manera modular y reutilizable. \\
\end{longtable}
